\documentclass{amsart}
% For dealing with references we use the comment environment
\usepackage{verbatim}
\newenvironment{reference}{\comment}{\endcomment}
%\newenvironment{reference}{}{}
\newenvironment{slogan}{\comment}{\endcomment}
\newenvironment{history}{\comment}{\endcomment}

% For commutative diagrams we use Xy-pic
\usepackage[all]{xy}

% We use 2cell for 2-commutative diagrams.
\xyoption{2cell}
\UseAllTwocells

% We use multicol for the list of chapters between chapters
\usepackage{multicol}

% This is generally recommended for better output
\usepackage{lmodern}
\usepackage[T1]{fontenc}

% For cross-file-references

% Package for hypertext links:
\usepackage{hyperref}

% For any local file, say "hello.tex" you want to link to please

% Theorem environments.
%
\theoremstyle{plain}
\newtheorem{theorem}[subsection]{Theorem}
\newtheorem{proposition}[subsection]{Proposition}
\newtheorem{lemma}[subsection]{Lemma}

\theoremstyle{definition}
\newtheorem{definition}[subsection]{Definition}
\newtheorem{example}[subsection]{Example}
\newtheorem{exercise}[subsection]{Exercise}
\newtheorem{situation}[subsection]{Situation}

\theoremstyle{remark}
\newtheorem{remark}[subsection]{Remark}
\newtheorem{remarks}[subsection]{Remarks}

\numberwithin{equation}{subsection}

% Macros
%
\def\lim{\mathop{\mathrm{lim}}\nolimits}
\def\colim{\mathop{\mathrm{colim}}\nolimits}
\def\Spec{\mathop{\mathrm{Spec}}}
\def\Hom{\mathop{\mathrm{Hom}}\nolimits}
\def\Ext{\mathop{\mathrm{Ext}}\nolimits}
\def\SheafHom{\mathop{\mathcal{H}\!\mathit{om}}\nolimits}
\def\SheafExt{\mathop{\mathcal{E}\!\mathit{xt}}\nolimits}
\def\Sch{\mathit{Sch}}
\def\Mor{\mathop{\mathrm{Mor}}\nolimits}
\def\Ob{\mathop{\mathrm{Ob}}\nolimits}
\def\Sh{\mathop{\mathit{Sh}}\nolimits}
\def\NL{\mathop{N\!L}\nolimits}
\def\CH{\mathop{\mathrm{CH}}\nolimits}
\def\proetale{{pro\text{-}\acute{e}tale}}
\def\etale{{\acute{e}tale}}
\def\QCoh{\mathit{QCoh}}
\def\Ker{\mathop{\mathrm{Ker}}}
\def\Im{\mathop{\mathrm{Im}}}
\def\Coker{\mathop{\mathrm{Coker}}}
\def\Coim{\mathop{\mathrm{Coim}}}

% Boxtimes
%
\DeclareMathSymbol{\boxtimes}{\mathbin}{AMSa}{"02}

%
% Macros for moduli stacks/spaces
%
\def\QCohstack{\mathcal{QC}\!\mathit{oh}}
\def\Cohstack{\mathcal{C}\!\mathit{oh}}
\def\Spacesstack{\mathcal{S}\!\mathit{paces}}
\def\Quotfunctor{\mathrm{Quot}}
\def\Hilbfunctor{\mathrm{Hilb}}
\def\Curvesstack{\mathcal{C}\!\mathit{urves}}
\def\Polarizedstack{\mathcal{P}\!\mathit{olarized}}
\def\Complexesstack{\mathcal{C}\!\mathit{omplexes}}
% \Pic is the operator that assigns to X its picard group, usage \Pic(X)
% \Picardstack_{X/B} denotes the Picard stack of X over B
% \Picardfunctor_{X/B} denotes the Picard functor of X over B
\def\Pic{\mathop{\mathrm{Pic}}\nolimits}
\def\Picardstack{\mathcal{P}\!\mathit{ic}}
\def\Picardfunctor{\mathrm{Pic}}
\def\Deformationcategory{\mathcal{D}\!\mathit{ef}}

\setlength{\emergencystretch}{3em}
\begin{document}
\title[Stacks Project, Tag 0B1Q]{461 --- General projective projections are locally immersions or etale at a fixed nonsingular point}
\maketitle
\noindent Copyright (C) 2005--2025 Johan de Jong. Stacks Project authors. Source: \href{https://stacks.math.columbia.edu/tag/0B1Q}{Tag 0B1Q}.

\noindent Editorial difficulty note (not author text). Difficulty H1: One must produce one nonempty parameter open avoiding both secant identifications and tangent collapse. The proof constructs two incidence schemes with strict dimension bounds, removes their constructible images, and then combines finiteness, tangent injectivity and the single-fibre criterion to obtain a local immersion or etaleness..

\noindent Editorial provenance note (not author text). History: excerpt from revision \texttt{a0444\allowbreak 6e57e\allowbreak c1fbc\allowbreak 252a8\allowbreak 71afc\allowbreak ec775\allowbreak 2fb28\allowbreak 07b14}, intersection.tex, lines 2257--2316. Reference presentation was adapted; original-author context and core support blocks were retained or added. The mathematical wording is unchanged. Licensed under GNU FDL 1.2 or later; no invariant sections or cover texts. See ../licenses/COPYING. Context blocks reproduce author definitions and supporting results; they are not additional counted problems.

\section{Projections}
\label{section-projection}

\noindent
Recall that we are working over a fixed algebraically closed ground field
$\mathbf{C}$. If $V$ is a finite dimensional vector space over $\mathbf{C}$
then we set
$$
\mathbf{P}(V) = \text{Proj}(\text{Sym}(V))
$$
where $\text{Sym}(V)$ is the symmetric algebra on $V$ over $\mathbf{C}$.
See Constructions, Example \href{https://stacks.math.columbia.edu/tag/0FCY}{example projective space (Tag 0FCY)}.
The normalization is chosen such that
$V = \Gamma(\mathbf{P}(V), \mathcal{O}_{\mathbf{P}(V)}(1))$.
Of course we have $\mathbf{P}(V) \cong \mathbf{P}^n_{\mathbf{C}}$ if
$\dim(V) = n + 1$. We note that $\mathbf{P}(V)$ is a nonsingular projective
variety.

\medskip\noindent
Let $p \in \mathbf{P}(V)$ be a closed point. The point $p$ corresponds to a
surjection $V \to L_p$ of vector spaces where $\dim(L_p) = 1$, see
Constructions, Lemma \href{https://stacks.math.columbia.edu/tag/01NA}{lemma apply (Tag 01NA)}.
Let us denote $W_p = \Ker(V \to L_p)$.
{\it Projection from $p$} is the morphism
$$
r_p : \mathbf{P}(V) \setminus \{p\} \longrightarrow \mathbf{P}(W_p)
$$
of Constructions, Lemma \href{https://stacks.math.columbia.edu/tag/01MY}{lemma morphism proj (Tag 01MY)}.
Here is a lemma to warm up.



\begin{lemma}
\label{lemma-projection-generically-immersion}
Let $V$ be a vector space of dimension $n + 1$.
Let $X \subset \mathbf{P}(V)$ be a closed subvariety.
Let $x \in X$ be a nonsingular point.
\begin{enumerate}
\item If $\dim(X) < n - 1$, then there is a nonempty Zariski open
$U \subset \mathbf{P}(V)$ such that for all closed points $p \in U$ the
morphism $r_p|_X : X \to r_p(X)$ is an
isomorphism over an open neighbourhood of $r_p(x)$.
\item If $\dim(X) = n - 1$, then there is a nonempty Zariski open
$U \subset \mathbf{P}(V)$ such that for all closed points $p \in U$ the
morphism $r_p|_X : X \to \mathbf{P}(W_p)$ is \'etale at $x$.
\end{enumerate}
\end{lemma}

\begin{proof}
Proof of (1). Note that if $x, y \in X$ have the same image under
$r_p$ then $p$ is on the line $\overline{xy}$.
Consider the finite type scheme
$$
T = \{(y, p) \mid
y \in X \setminus \{x\},\ p \in \mathbf{P}(V),\ p \in \overline{xy}\}
$$
and the morphisms $T \to X$ and $T \to \mathbf{P}(V)$ given by
$(y, p) \mapsto y$ and $(y, p) \mapsto p$.
Since each fibre of $T \to X$ is a line, we see that
the dimension of $T$ is $\dim(X) + 1 < \dim(\mathbf{P}(V))$.
Hence $T \to \mathbf{P}(V)$ is not surjective. On the other hand,
consider the finite type scheme
$$
T' = \{p \mid
p \in \mathbf{P}(V) \setminus \{x\},
\ \overline{xp}\text{ tangent to }X\text{ at }x\}
$$
Then the dimension of $T'$ is $\dim(X) < \dim(\mathbf{P}(V))$.
Thus the morphism $T' \to \mathbf{P}(V)$ is not surjective either.
Let $U \subset \mathbf{P}(V) \setminus X$ be nonempty open and disjoint
from these images; such a $U$ exists because the images of $T$ and $T'$
in $\mathbf{P}(V)$ are constructible by
Morphisms, Lemma \href{https://stacks.math.columbia.edu/tag/054J}{lemma chevalley (Tag 054J)}.
Then for $p \in U$ closed the projection
$r_p|_X : X \to \mathbf{P}(W_p)$ is injective on the
tangent space at $x$ and $r_p^{-1}(\{r_p(x)\}) = \{x\}$.
This means that $r_p$ is unramified at $x$
(Varieties, Lemma \href{https://stacks.math.columbia.edu/tag/0B2G}{lemma injective tangent spaces unramified (Tag 0B2G)}),
finite by Lemma \href{https://stacks.math.columbia.edu/tag/0B1P}{lemma projection generically finite (Tag 0B1P)},
and $r_p^{-1}(\{r_p(x)\}) = \{x\}$ thus \'Etale Morphisms, Lemma
\href{https://stacks.math.columbia.edu/tag/04DG}{lemma finite unramified one point (Tag 04DG)} applies and
there is an open neighbourhood $R$ of $r_p(x)$
in $\mathbf{P}(W_p)$ such that $(r_p|_X)^{-1}(R) \to R$ is a
closed immersion which proves (1).

\medskip\noindent
Proof of (2). In this case we still conclude that the morphism
$T' \to \mathbf{P}(V)$ is not surjective.
Arguing as above we conclude that for $U \subset \mathbf{P}(V)$
avoiding $X$ and the image of $T'$, the projection
$r_p|_X : X \to \mathbf{P}(W_p)$ is \'etale at $x$ and finite.
\end{proof}
\end{document}
